\documentclass[arxiv,usenatbib]{iupartex}
% geometry, xcolor, graphicx, amsmath, amssymb, hyperref, natbib and dcolumn
% latex packages are used in the class file, do not redeclare them.

%%%%%%%%%%%%%% Attention to authors %%%%%%%%%%%%%%%%%%%%%%
% template has two output options, referee and arxiv
%%%%%%%%%%%%%%%%%%%%%%%%%%%%%%%%%%%%%%%%%%%%%%%%%%%%%%%%%%
% Don't change following lines
\usepackage{newtxtext,newtxmath}
\usepackage[T1]{fontenc}

%%%%% Authors - Place Your Own Packages Here %%%%%
% Only include extra packages if you really need them. Common packages are:
%\usepackage{subcaption}
%\setcitestyle{square,numbers}

%%%% your own commands here %%%%%%%%%%%%%%%
\newcommand{\textbit}[1]{\textbf{\textit{#1}}}
\newcommand{\inspace}[1]{\vspace*{#1}\noindent}

%%%%%%%%%%%%%%%%%%% Title Page %%%%%%%%%%%%%%%%%%%
%\title[Running title]{article title}
\title[PAR LATEX template]{PAR LATEX template - title goes here}

% The list of authors, and the short list which is used in the headers.
% If you need two or more lines of authors, add an extra line using \newauthor
%% for corresponding author put \cc after superscripts number
\author[Surname et. al]{%
E. A. Surname$^{1\cc}$,\orcid{0000-000X-XXXX-AAAA}
A. N. Other$^{2}$,\orcid{0000-000X-XXXX-BBBB}
Third Author $^{2,3}$\orcid{0000-000X-XXXX-CCCC}
and
Fourth Author $^{3}$\orcid{0000-000X-XXXX-DDDD}
% do not delete the following line
\affsep \\
% List of institutions with line breaks
$^1$Istanbul University, Faculty of Science, Department of Physics, Vezneciler, 34134, Istanbul, T\"{u}rkiye\\
$^2$Department, Institution, Street Address, City Postal Code, Country\\
$^3$Another Department, Another Institution, Street Address, City Postal Code, Country
}

%corresponding author name and email
\corres{E. A. Surname}{author@email}

% Enter the current year, for the copyright statements etc.
\pubyear{20XX}

%%%%%%%% Authors are not allowed to edit %%%%%%%%%
% Don't change the lines up to \maketitle
\doiheader{XXXXXXX/PAR.20XX.00000}
\date{
	\pSubmit{00.00.0000} 
	\pRevReq{00.00.0000}
	\pLastRevRec{00.00.0000}
	\pAccept{00.00.0000}
	\pPubOnl{00.00.0000}
}
\volume{0}
\volnumber{1}

\begin{document}
\label{firstpage}
\pagerange{\pageref*{firstpage}--\pageref*{lastpage}}
\maketitle

% Abstract of the paper
\begin{abstract}
This is a simple template for authors to write new PAR papers. The abstract should briefly describe the aims, methods, and main results of the paper. It should be a single paragraph not more than 300 words. No references should appear in the abstract.
\end{abstract}

% Select between one and six entries from the list of approved keywords.
% Don't make up new ones.
\begin{keywords}
keyword1 -- keyword2 -- keyword3
\end{keywords}

%%%%%%%%%%%%%%%%%%%%%%%%%%%%%%%%%%%%%%%%%%%%%%%%%%

%%%%%%%%%%%%%%%%% BODY OF PAPER %%%%%%%%%%%%%%%%%%

\section{Introduction}
This is a simple template for authors to write new PAR papers. See par\_sample.tex for a more complex example, and par\_guide.tex for a full user guide. 

All papers should start with an Introduction section, which sets the work in context, cites relevant earlier studies in the field by \cite{aut2018}, and describes the problem the authors aim to solve.
 

\section{Section}
Manuscripts should be divided into ordered and numbered sections. These sections should contain sufficient details to allow others to replicate and build on published results. New methods and protocols should be described in detail while well-established methods can be briefly described and appropriately cited. All headings (except Introduction and Conclusion) can be divided into subsections. 

\subsection{SubSection}
\subsubsection{Sub-SubSection}
\inspace{2mm}%
Bulleted lists look like this:
\begin{itemize}
	\item First bullet; 
	\item Second bullet;
	\item Third bullet.
\end{itemize}

\noindent Numbered lists look like this:
\begin{enumerate}
	\item First item
	\item Second item
	\item Third item
\end{enumerate}

\noindent Authors should provide a concise and precise description of the theoretical/experimental results, their interpretation as well as the conclusions that can be drawn in these sections.

\inspace{6mm}%
\textbit{Figures, Tables and Schemes}

\inspace{2mm}%
All figures and tables should be cited in the main text as Figure 1, Table 1, etc.

\begin{figure}
	\begin{center}
		\label{fig1}
		\includegraphics[width=0.92\columnwidth]{image1.png}
		\caption{This is a figure. Schemes follow the same formatting. If there are multiple panels, they should be listed as: (a) Description of what is contained in the first panel. (b) Description of what is contained in the second panel. Figures should be placed in the main text near to the first time they are cited.} 
	\end{center}
\end{figure}

\setlength{\tabcolsep}{5mm} % adjust space between the table columns
\renewcommand{\arraystretch}{1.0} % adjust space between the table rows
\begin{table}
	\label{table1}
	\caption{This is a table caption. Tables should be placed in the main text near to the first time they are cited}
	\centering
	\begin{tabular}{lll}
		\hline
		Title 1    & Title 2   & Title 3 \\
		\hline
		Data 1     & Data 3$^1$    & Data 5 \\
		Data 2     & Data 4    & Data 6 \\
		\hline
		\multicolumn{3}{@{}l}{\footnotesize{$^1$Tables may have footer}}
	\end{tabular}
\end{table}%

\inspace{2mm}%
\textbf{NOTE:} Tables, if preferred, can also be prepared in landscape format.

\inspace{6mm}%
\textbit{Mathematics}

\noindent Scalar variables are \textit{italic}; vectors are \textbit{bold italic} (no arrows); matrices are bold; dot products are denoted by a bold centred dot $\bullet$, cross-products by a bold multiplication sign $\times$. Differential d, complex i, exponential e, sin, cos, tan, log, etc., are not italic. Sub/superscripts that are physical variables are italic, while those that are merely labels are roman (e.g. $C_t$ and $F_v$ but $T_{\rm eff}$ and $b_{\rm max}$). 

\inspace{7mm}%
\textbit{Equations}

\inspace{3mm}%
This is the example of a simple equation:
\begin{equation}
	x=1
\end{equation}

\inspace{3mm}
This is the example of a longer equation:
\begin{equation}
	y=a+b+c+d+e+f+g+h+i+j+k+l+m+n+o+p
\end{equation}

\inspace{3mm}%
This is the example of a sequence of equations:
\begin{equation}
	\begin{aligned}
		x &= 1 \\
		y &= 2 \\
		z &= 3
	\end{aligned}
\end{equation}

\section{Conclusion}
This last section should briefly summarise what has been done, and describe the final conclusions which the authors draw from their work.

\section*{Acknowledgements}
Here you can thank helpful colleagues, acknowledge funding agencies and facilities used. It should be kept short.

%%%%%%%%%%%%%%%%%%%% REFERENCES %%%%%%%%%%%%%%%%%%
% The best way to enter references is to use BibTeX:

%\bibliographystyle{mnras}
%\bibliography{example} % if your bibtex file is called example.bib

%\bibliographystyle{ieeer}
\begin{thebibliography}{}
\bibitem[Author1 et al. (2018)]{aut2018} 
Author 1, S., Author 2, T., Author 3, U., 2018, Journal Abbreviation, 10, 142-149. 

\bibitem[Author1 (2007)]{aut2007}
Author 1, L., 2007, The title of the contribution. The Book Title; Editor 1, F., Editor 2, A., Eds.; Publishing House: City, Country, 32–58.

\bibitem[Author1 and Author2 (2008)]{aut2008} 
Author 1, A., Author 2, B., 2008, Book Title, 3rd ed.; Publisher: Publisher Location, Country, 154–196.

\bibitem[Author1 and Author2 (2008)]{aut2008a} 
Author 1, A.B., Author 2, C., 2008, Abbreviated Journal Name year, phrase indicating stage of publication (submitted; accepted; in press).

\bibitem[Author1 (2015)]{aut2015} 
Author 1, A. B., 2015, Title of Thesis. Level of Thesis, Degree-Granting University, Location of University, 

\bibitem[Title of Site (2000)]{weburl}
Title of Site. Available online: URL (accessed on Day Month Year).

\bibitem[Author1 and Author2 (2012)]{aut2012}
Author 1, A.B.; Author 2, C. Title of Unpublished Work. Abbreviated Journal Name year, phrase indicating stage of publication (submitted; accepted; in press).

\bibitem[Author1 et al. (2018)]{aut2018a}
Author 1, T., Author 2, S., Author 3, K., et al., 2018, preprint (arXiv: 0245.35874) 

\end{thebibliography}


\section*{Supplementary}
If you want to present additional material which would interrupt the flow of the main paper, it can be placed in an Appendix which appears after the list of references.

% Don't change these lines
\bsp	% typesetting comment
\label{lastpage}
\end{document}